\section{Modele językowe}
W części eksperymentalnej wykorzystano modele Llama-3.1-8B-Instruct~\cite{grattafiori2024llama3herdmodels} oraz Mistral-7B-Instruct-v0.3~\cite{jiang2023mistral7b}. Są to modele otwartoźródłowe z upublicznionymi paramaterami wewnętrznymi. Było to ważne kryterium ze względu na konieczność dostępu do tych informacji przy okazji implementacji metody obrony opartej na inżynierii aktywacji, a także wymagane były przy białoskrzynkowym ataku GCG.
Zostały one dostrojone podążać za poleceniami użytkownika i odpowiadać na zadane pytania, czego najczęstszym celem są niebezpieczne prompty.
W literaturze do badań często wybierano właśnie modele dostarczane przez firmę Meta, lecz z rodziny Llama~2. W tej pracy uznano, że wybrana rodzina tych modeli stanowi bardziej adekwatny punkt odniesienia dla współczesnej rzeczywistości badań nad bezpieczeństwem LLM.
Wybranie w badaniu również modelu Mistral pozwoliło zminimalizować ryzyko sformułowania błędnych wniosków na temat skuteczności metod obrony lub ataku, które wynikać by mogły ze specyfiki pojedynczego modelu lub jego zbioru uczącego.
Pozwoliło to sprawdzić, czy badana podatność jest również obecna w innym modelu oraz umożliwiło ocenę stosowalności metody obrony niezależnie od dostawcy danego LLM. Ponadto, mają one podobną liczbę parametrów, mianowicie 8 i 7 miliardów, co umożliwiło porównanie rezultatów na tej samej maszynie.
Wybór modeli o takim rozmiarze miał również znaczenie przy selekcji ze względu na dostępne zasoby obliczeniowe.
